\documentclass[12pt]{article}

%% --- Packages ---
\usepackage[margin=1in]{geometry}
\usepackage{setspace}
\doublespacing
\usepackage{amsmath,amssymb}
\usepackage{booktabs}
\usepackage{caption}
\usepackage{graphicx}
\usepackage{hyperref}
\usepackage[authoryear,round]{natbib}
\usepackage{threeparttable}
\usepackage{rotating}
\usepackage{lscape}
\usepackage{array}
\usepackage{xcolor}
\usepackage[T1]{fontenc}

\hypersetup{colorlinks=true, linkcolor=black, citecolor=black, urlcolor=black}

%% --- Title ---
\title{\large\textbf{Wildfire Incidence, Fire Severity, and Local Public Finance: \\
Evidence from Western U.S.\ County Governments}}

\author{
  Yong Chen\thanks{[Affiliation]. Email: [email]. I thank [acknowledgments].}
}

\date{June 2026 \\ \smallskip \small Preliminary draft. Comments welcome.}

\begin{document}
\maketitle

\begin{abstract}
\noindent I estimate the causal effect of wildfire exposure on county government fiscal
outcomes along two margins of the wildfire shock: the \textit{extensive margin}
(fire incidence---whether a county experiences any qualifying wildfire) and the
\textit{intensive margin} (fire severity---whether the fire caused detectable
stand-replacement burning). Using quinquennial Census of Governments data from
2002 to 2022 for 344 county governments across eight western states, I implement
two complementary staggered difference-in-differences designs with Callaway and
Sant'Anna (2021) doubly-robust estimation, exploiting the 2012 USFS Wildfire
Potential raster as a predetermined matching variable. In the extensive-margin
specification, 293 counties experienced a first qualifying wildfire
($\geq\!1{,}000$ acres) between 2013 and 2021; the 51 never-fired counties serve
as controls (ESS: 11.1 after reweighting). In the intensive-margin specification,
246 counties experienced a first high-severity wildfire (MTBS High-severity
dNBR class detected); the 98 counties without a qualifying high-severity
fire---including those with only low-to-moderate-severity fires---serve as
controls (ESS: 19.6). Both margins yield statistically insignificant fiscal
effects. The extensive-margin ATT for fiscal balance is \$\!-649 per capita
(SE\,=\,525; $p = 0.217$); the intensive-margin ATT is \$\!-278 (SE\,=\,245;
$p = 0.258$). Capital expenditures show the most consistent signal: \$\!-226
(extensive) and \$\!-181 (intensive) per capita. For the intensive-margin
specification, all four heterogeneity-robust estimators (Callaway-Sant'Anna DR,
Sun-Abraham, Synthetic DiD, two-stage DiD) agree on negative signs---a
cross-estimator consensus attributable to the closer counterfactual provided
by the expanded low-severity-fire control group.
\medskip

\noindent\textbf{Keywords:} wildfire; fire severity; local public finance; fiscal federalism;
difference-in-differences; western United States

\noindent\textbf{JEL:} H72, Q54, H71, R11
\end{abstract}

\newpage

%% ============================================================
\section{Introduction}
%% ============================================================

Western wildfires have grown in frequency, severity, and geographic reach over the
past two decades. Between 2013 and 2021---the treatment window in this study---293
of 344 county governments across eight Pacific Coast and Rocky Mountain states
experienced at least one qualifying wildfire ($\geq\!1{,}000$ acres). Among these,
246 counties experienced a fire that the USFS Monitoring Trends in Burn Severity
(MTBS) program classified as producing a distinct high-severity burning class---fires
with stand-replacement behavior and severe vegetation mortality. Beyond their direct
costs to households and ecosystems, large wildfires impose a first-order fiscal
challenge on the county governments responsible for emergency response,
infrastructure repair, and baseline public services: property tax revenues may fall
as fire damages assessed values, while expenditure demands rise simultaneously.
How much fiscal stress wildfires impose---and whether severity amplifies the effect
beyond mere incidence---is a central concern for state and local fiscal policy,
yet causal evidence for county governments outside California is absent.

This paper provides that evidence along two distinct margins. The
\textit{extensive margin} asks whether experiencing any qualifying wildfire
($\geq\!1{,}000$ acres) affects county fiscal outcomes; the 293 fire-affected
counties are compared to 51 counties that experienced no qualifying fire.
The \textit{intensive margin} asks whether fire \textit{severity}---holding
fire occurrence approximately fixed---amplifies the fiscal effect; the 246
counties with high-severity fires are compared to 98 counties without a
qualifying high-severity fire (including 47 counties with only
low-to-moderate-severity fires). I use quinquennial Census of Governments data
from 2002 to 2022, restricting to county governments in California, Colorado,
Idaho, Montana, Oregon, Utah, Washington, and Wyoming---the eight states with
the highest wildfire incidence and the greatest share of high-hazard county
governments. The panel covers 344 counties across five census years, yielding
1,720 observations.

Identification in both designs exploits the predetermined 2012 USFS Wildfire
Potential (WFP) raster as a matching variable in a staggered
difference-in-differences (DiD) framework. The primary estimator is Callaway and
Sant'Anna's \citeyearpar{callaway2021} doubly-robust group-time estimator.
Within a given wildfire hazard class, the timing of a county's first large
fire in 2013--2021 depends on stochastic weather shocks---drought, lightning,
and wind---that are plausibly orthogonal to fiscal trends conditional on
baseline hazard and demographic covariates. I additionally verify pre-treatment
parallel trends for both designs. As robustness checks for the intensive-margin
specification I implement Sun and Abraham's \citeyearpar{sun2021} TWFE estimator,
Arkhangelsky et al.'s \citeyearpar{arkhangelsky2021} Synthetic DiD, and
Gardner's \citeyearpar{gardner2022} two-stage DiD.

The main finding is that \textit{neither margin generates statistically significant
fiscal effects.} For the extensive margin, the doubly-robust ATT for fiscal balance
is \$\!-649 per capita (SE\,=\,525; $p = 0.217$; 95\% CI: [\$\!-1{,}678, \$381]).
For the intensive margin, the ATT is \$\!-278 (SE\,=\,245; $p = 0.258$; 95\% CI:
[\$\!-758, \$203]). Capital expenditures show the most consistent signal: \$\!-226
per capita (extensive, $p = 0.166$) and \$\!-181 per capita (intensive, $p = 0.142$),
suggesting that whatever fiscal adjustment occurs operates primarily through deferred
capital investment. For the intensive-margin specification, all four estimators agree
on negative signs for fiscal balance---ranging from \$\!-126 (Sun-Abraham) to
\$\!-278 (C\&S DR)---a cross-estimator consensus that does not hold in the
extensive-margin specification, where covariate adjustment reverses the sign
relative to unadjusted estimators. The imprecision reflects the limited quinquennial
observation window: effective control-group sizes after doubly-robust reweighting
are 11.1 (extensive) and 19.6 (intensive) counties.

This paper makes three contributions. First, it provides causal estimates of wildfire fiscal effects for county
governments across eight western states, extending \citet{liao2022} from
California municipalities to county governments with diverse property tax
institutions, evaluated at both the extensive (incidence) and intensive
(severity) margins; a California-excluded robustness check yields qualitatively
identical null findings (Table~\ref{tab:robustness}, row~5), confirming that
results are not driven by California's Proposition~13 environment.
Second, it demonstrates that defining treatment by fire severity rather than
fire incidence improves the identification environment: the expanded
low-to-moderate-severity control pool reduces observable selection and
yields cross-estimator sign agreement. Third, the null result, while
consistent with the range of plausible effects, reflects a power constraint
that annual CoG data would resolve: Table~\ref{tab:mde} shows that the minimum
detectable effect is \$687--\$1{,}471 per capita for fiscal balance, well above the
\citet{liao2022} California benchmark of \$97 per capita.

The paper proceeds as follows. Section~2 reviews the related literature.
Section~3 presents the data and both treatment definitions. Section~4 details
the empirical strategy. Section~5 reports main results for both margins.
Section~6 presents robustness checks. Section~7 concludes.

%% ============================================================
\section{Related Literature}
%% ============================================================

\textbf{Natural disasters and local public finance.} The disaster-finance literature
has established that natural disasters impose persistent fiscal costs on local
governments, but the evidence base rests almost entirely on hurricanes and floods.
\citet{deryugina2017} showed that post-hurricane government transfers to affected
counties exceed direct disaster aid in present value, producing approximate net
fiscal neutrality; the mechanism is federal social insurance absorption rather
than own-source revenue recovery. \citet{jerch2023} extended this framework to
municipal governments, finding that major hurricanes reduced local tax revenues and
raised debt costs persistently for over a decade, with effects concentrated in
minority-majority municipalities. Neither paper examines wildfire, where the damage
mechanism---forest land and vegetation rather than residential structure stock---the
federal aid architecture (FEMA declarations are less automatic), and the affected
geography (rural western counties with thin property tax bases) differ substantially
from hurricane-exposed Atlantic and Gulf Coast governments.

\textbf{Wildfire economics.} Causal research on wildfires has concentrated on
property values and residential development \citep{baylis2023}, health outcomes
from smoke exposure \citep{borgschulte2024}, and insurance markets.
\citet{liao2022} is the sole paper providing causal estimates of wildfire fiscal
impacts: using quasi-random fire perimeter exposure within California, they found
that municipalities experienced a 25-percentage-point increase in the probability
of running a budget deficit post-fire, with expenditures growing 17.3\% against
revenues growing 10.5\%, generating a net fiscal loss of approximately \$97 per
capita. Their setting, however, is California-specific in two important ways.
First, Proposition 13 creates a unique property tax dynamic: assessed values are
frozen at purchase price and only reset upon sale, which means post-fire property
sales generate reassessment events that temporarily raise assessed values and
transfer tax revenues even as fire-damaged properties lose market value. Second,
their sample is municipal governments, which operate under different revenue
authorities than county governments and are concentrated in the high-density,
high-value coastal and suburban areas of California rather than in the rural
forested communities most exposed to wildfire in other western states.
The proposed paper directly extends \citet{liao2022} to county governments
across eight states without Proposition 13.

\textbf{Property tax dynamics.} Property tax revenues do not immediately reflect
market value changes. \citet{lutz2011} documented a four-year lag between market
value peaks and property tax revenue troughs during the 2008 housing crisis.
\citet{chen2025} show that a 1\% market value change translates to less than 0.30\%
change in assessed value within three years, with wide cross-jurisdiction variance
driven by assessor practices and state law. This assessment lag implies that the
most visible property tax revenue response to wildfire will emerge two to four years
after fire occurrence---the horizon captured by the 2017 and 2022 CoG census
years for the treatment cohort entering in 2013--2016.

\textbf{Staggered DiD methodology.} Standard two-way fixed-effects (TWFE) estimators
are biased when treatment effects are heterogeneous across cohorts and calendar time
\citep{baker2022, sun2021}. I use the Callaway and Sant'Anna \citeyearpar{callaway2021}
group-time estimator, which identifies cohort-specific average treatment effects and
aggregates them with user-specified weights, as the primary specification.
I report \citet{sun2021} and \citet{arkhangelsky2021} estimates as robustness checks.
Pre-trend tests follow \citet{roth2022}, who cautions against mechanically accepting
pre-treatment parallel trends as sufficient for identification.

%% ============================================================
\section{Data}
%% ============================================================

\subsection{Sample Construction}

The sample comprises county governments in eight western states: California, Colorado,
Idaho, Montana, Oregon, Utah, Washington, and Wyoming. I restrict to counties
present in the Census of Governments quinquennial census for all five census years
(2002, 2007, 2012, 2017, 2022), yielding a balanced panel of 344 counties and 1,720
county-year observations. Two Montana consolidated city-counties (Deer Lodge and
Silver Bow) are excluded because they are not classified as county governments in
the CoG entity type system.

I implement two treatment definitions that correspond to the extensive and intensive
margins of wildfire exposure.

\textit{Extensive margin (fire incidence).} Treatment is a county's first qualifying
MTBS wildfire ($\geq\!1{,}000$ acres) within the 2013--2021 window, regardless of
severity. Under this definition, 293 counties are treated: 208 experienced a first
qualifying fire in 2013--2016 (early cohort, g\,=\,2017; observed post-treatment in
the 2017 and 2022 censuses) and 85 experienced a first qualifying fire in 2017--2021
(late cohort, g\,=\,2022; observed in the 2022 census). The 51 counties with no
qualifying MTBS fire in 2013--2021 form the never-treated control group
(effective sample size: 11.1 counties after doubly-robust reweighting).

\textit{Intensive margin (fire severity).} Treatment is a county's first
\textit{high-severity} qualifying wildfire within the 2013--2021 window. A fire
qualifies as high-severity if the MTBS program identified a distinct High-severity
dNBR classification class, as indicated by the fire-specific \texttt{HIGH\_THRES}
attribute (between 100 and 5,000 dNBR units) in the MTBS perimeter file. A
\texttt{HIGH\_THRES} value in this range indicates that MTBS analysts fitted a
fire-specific dNBR histogram with a classifiable High-severity tail, corresponding
to stand-replacement burning with severe vegetation mortality. Fires with
\texttt{HIGH\_THRES}\,=\,9,999 (sentinel: no distinct high-severity class) are
classified as low-to-moderate-severity. Under this definition, 246 counties are
treated: 173 in the early cohort (g\,=\,2017) and 73 in the late cohort (g\,=\,2022).
The 98 counties without a qualifying high-severity fire in 2013--2021---including
47 counties that experienced only low-to-moderate-severity fires---form the
never-treated control group (effective sample size: 19.6 counties after
doubly-robust reweighting).

The intensive-margin control group is strictly larger than the extensive-margin
control group: it includes the same 51 zero-fire counties plus 47 counties with
low-to-moderate-severity fires. These 47 counties provide an expanded, closer
counterfactual for counties experiencing high-severity fires, because they share
similar wildfire hazard profiles and pre-treatment fire histories with the treated
group while differing primarily on fire severity.

\subsection{Fiscal Outcomes}

Fiscal data come from the Census of Governments quinquennial census, which provides
complete coverage of all government units. For the 2017 and 2022 censuses, which
use the Individual Unit File format (a 32-character fixed-width record), aggregate
summary codes (total general revenue, total general expenditure) are absent.
I reconstruct these aggregates from line-item codes following the published
itemization guide, summing across federal and state intergovernmental revenue
(B-codes and C-codes), own-source revenue (T, U, and A codes), and direct general
expenditure (E, F, G, and J codes).

All fiscal outcomes are expressed per capita in 2019 dollars, using ACS 5-year
estimates for population denominators and the national CPI-U for deflation.
Fiscal year data are assigned to the calendar year in which the fiscal year begins:
a fiscal year ending June 30, 2017 is coded to 2016. In the eight states in the
sample, all county fiscal years end on June 30, except Idaho (September 30);
the resulting treatment cohort assignments are unaffected because the timing
difference is less than one calendar year.

The headline outcome is \textit{fiscal balance per capita} (total general revenues
minus total general expenditures), which captures the net fiscal burden in a
single directly interpretable unit and maps to the \citet{liao2022} budget deficit
probability result. Secondary outcomes decompose the fiscal balance into revenue
(total revenue, property tax revenue, intergovernmental transfers) and
expenditure (total expenditures, capital outlays) components.

\subsection{Wildfire Hazard and Treatment}

The primary matching variable is the 2012 USFS Wildfire Potential (WFP) raster
\citep[RDS-2015-0045]{scott2013}, which provides a continuous
index of wildfire likelihood and intensity at 270m resolution. The 2012 vintage is
predetermined for all fires from 2013 onward: it was finalized before the 2013 fire
season and reflects FSim fire modeling inputs calibrated on data through 2011.
County-level WFP is the area-weighted mean of raster cells intersecting the county
boundary, computed in the Albers Equal Area projection (EPSG:5070). WFP values
in the sample range from near-zero (urban or desert counties) to over 1,000 in
high-hazard forested counties.

Fire perimeters, treatment assignment, and severity attributes come from the
Monitoring Trends in Burn Severity (MTBS) database, which provides satellite-derived
fire perimeters for all fires exceeding 1,000 acres in the western US from 1984
onward. Each MTBS fire record includes the \texttt{HIGH\_THRES} attribute: the
fire-specific dNBR threshold above which pixels are classified as High-severity.
This threshold is set separately for each fire using the fire's dNBR histogram;
a value of 9,999 indicates the histogram lacked a classifiable High-severity tail,
meaning high-severity burning either did not occur or could not be isolated.
I intersect fire perimeters with county boundaries using GeoPandas and assign each
county the year of its first qualifying \textit{high-severity} fire in 2013--2021.
Counties with no qualifying high-severity fire in this window are never-treated.
I construct a 100 km smoke exclusion buffer around all fire perimeters. For fiscal
outcomes, which operate through property damage and local spending channels rather
than the smoke inhalation channel, the buffer restriction is not applied in the
main analysis.

\subsection{Covariates}

Propensity score estimation uses: WFP quintile (1--5), an indicator for any
qualifying fire before 2013, log pre-2013 cumulative acres burned, USDA Rural-Urban
Continuum Code (2013 vintage), log median household income, poverty rate,
uninsurance rate, and log population density---all measured at the 2012 CoG census
or from the 2007--2011 ACS 5-year estimates. Pre-treatment fiscal baselines
(log revenue per capita, log property tax per capita, log long-term debt per capita
in 2012) are included in the propensity score to account for pre-existing fiscal
differences between counties that will and will not experience wildfire.

\subsection{Summary Statistics}

Table~\ref{tab:balance} reports pre-treatment means and standardized mean differences
(SMD) for treated (high-severity fire) and control (no high-severity fire) counties,
before and after propensity-score inverse-probability weighting. In the unweighted
comparison, treated counties have substantially higher wildfire potential (WFP mean
1,064 vs.\ 208; SMD = 0.89), higher rates of pre-2013 fire history (92\% vs.\ 56\%),
and higher pre-2013 fire acres burned (SMD = 1.20). After IPW, WFP differences
compress to SMD = 0.13, and pre-2013 fire history is well-balanced (SMD = $-$0.03).
The expanded control group of 98 counties---enriched by counties with low-to-moderate
severity fires---provides a more relevant comparison set and improves the effective
sample size (ESS) of the reweighted control group to 19.6 (from 11.1 in an all-fire
analysis), a 77\% improvement in effective statistical power.

%% ============================================================
\section{Empirical Strategy}
%% ============================================================

\subsection{Estimating Equation}

The primary estimator is the group-time average treatment effect (ATT) of
\citet{callaway2021}. Let $g \in \{2017, 2022\}$ denote the first calendar year
in which county $i$'s treatment cohort first appears post-treatment in a CoG
census year, and let $\infty$ denote never-treated. The group-time ATT is:
%
\begin{equation}
  ATT(g, t) \;=\; \mathbb{E}\!\left[Y_{it}(g) - Y_{it}(\infty) \mid G_i = g\right]
  \label{eq:att_gt}
\end{equation}
%
\noindent where $Y_{it}(g)$ is the potential outcome for county $i$ in census year
$t$ under first treatment in cohort $g$, and $Y_{it}(\infty)$ is the never-treated
potential outcome. I estimate equation~(\ref{eq:att_gt}) using the doubly-robust
(DR) estimator of \citet{callaway2021}, which combines propensity-score and
outcome-regression adjustments so that the estimate is consistent if either
the propensity score or the outcome model is correctly specified.

The propensity score model is a logistic regression of ever-treated status on
the covariates described in Section~3. The outcome regression includes the same
covariate vector. I use \texttt{base\_period = ``varying''}, which sets a
pre-treatment comparison period separately for each group-time cell, and
\texttt{control\_group = ``nevertreated''}, restricting comparisons to the 98
never-treated counties (those without a qualifying high-severity fire in 2013--2021).
Standard errors are clustered at the county level and computed via block bootstrap
with 999 replications.

The primary reported estimate is the simple aggregated ATT:
%
\begin{equation}
  \widehat{ATT} \;=\; \sum_{g,t:\,t \geq g} \hat{\omega}(g,t) \cdot \widehat{ATT}(g,t)
  \label{eq:att_simple}
\end{equation}
%
\noindent where $\hat{\omega}(g,t) \propto \Pr(G_i = g)$ weights each cohort-period
cell by cohort size. I also report cohort-specific aggregations (ATT for
g\,=\,2017 and g\,=\,2022 separately) and event-study dynamic aggregations
(ATT at relative time $\ell = t - g$ for $\ell \in \{-10, -5, 0, 5\}$ in
quinquennial CoG years).

\subsection{Identifying Assumption}

The key identifying assumption is conditional parallel trends: conditional on the
WFP quintile, pre-treatment fire history, and other covariates, treated and
never-treated counties would have followed the same fiscal trajectory absent
wildfire. This assumption is most plausible in the DR framework, which adjusts
for observable selection into treatment. The fire literature provides a natural
justification: within a given wildfire hazard class, the timing of a county's
first large fire in 2013--2021 depends on stochastic weather shocks---drought
conditions, lightning ignition probability, and wind speed---that are plausibly
orthogonal to the county's fiscal trend conditional on baseline hazard level.

\subsection{Pre-trends and Power}

I test for pre-treatment parallel trends using the dynamic ATT estimates at
$\ell = -10$ and $\ell = -5$ (census years 2002 and 2007 for the g\,=\,2017 cohort).
A joint chi-squared test of the null that all pre-treatment dynamic ATTs equal
zero is reported for each outcome. Pre-trend tests pass comfortably for all
outcomes (minimum $p = 0.50$ for fiscal balance; minimum $p = 0.60$ across all
headline outcomes), consistent with the conditional parallel trends assumption.
The strong pre-trend performance reflects the improved comparability of the
expanded control group: counties with low-to-moderate-severity fires share
similar WFP profiles and pre-treatment fire histories with high-severity-fire
counties, making pre-treatment fiscal trajectories more parallel.

\citet{roth2022} cautions that passing pre-trend tests does not validate the parallel
trends assumption, because the test lacks power to detect confounding trends when
pre-treatment variation is limited. In this application, the three quinquennial
pre-treatment observations (2002, 2007, 2012) provide limited pre-trend signal.
Quantifying this limitation: for fiscal balance under the intensive margin,
the maximum pre-trend slope that passes the test undetected (1.96 times the
maximum pre-period bootstrap SE, a conservative scalar bound from
\citealt{roth2022}) is approximately \$795 per capita per quinquennium---nearly
three times the ATT point estimate of \$\!-278. For capital expenditures the
bound is \$289 per quinquennium (1.6 times the ATT). The pre-trend point
estimates are small (fiscal balance at $\ell = -15, -10, -5$: \$58, \$86,
\$237; none significant), providing qualitative reassurance, but the design
cannot rule out moderate confounding trends. Pre-trend results should be read
as consistent with, not confirmation of, conditional parallel trends.

\subsection{Threats to Identification}

\textbf{Property assessment lag.} Property tax revenues respond to wildfire with a
multi-year lag: assessors in most western states reassess on two- to three-year
cycles, and California's Proposition 8 calamity reassessment provisions create
a California-specific mechanism that allows immediate reassessment at owner request.
The quinquennial CoG census captures fiscal outcomes roughly four to nine years
post-fire for g\,=\,2017 counties and one to five years post-fire for g\,=\,2022
counties, ensuring that medium-run tax revenue responses are captured. The
g\,=\,2022 cohort, observed in only one post-treatment census year, provides
limited power for this cohort.

\textbf{FEMA endogeneity.} Federal disaster declarations are endogenously triggered by
fire severity, creating a mechanical correlation between fire damage and
intergovernmental revenue through FEMA Hazard Mitigation and Individual Assistance
programs. Intergovernmental revenue is therefore treated as a mechanism outcome
rather than a control variable. The main specification for fiscal balance uses
total revenue net of all intergovernmental transfers, isolating own-source
fiscal capacity.

\textbf{Propensity score specification.} The propensity score is a logistic
regression of ever-treated status on the covariate vector $\mathbf{X}$.
Separate propensity score models are estimated for the extensive margin
(ever-treated = any qualifying fire) and the intensive margin (ever-treated =
any high-severity fire). The C\&S doubly-robust estimator uses these
margin-specific PS weights together with an outcome regression to construct
the ATT. Standard errors are obtained by block bootstrap over counties.

\textbf{Spatial autocorrelation.} County clustering of standard errors accounts
for within-county serial correlation but not for cross-county spatial correlation.
Adjacent western counties share watershed boundaries and face common economic
shocks (regional drought, statewide fiscal rules) that may correlate their fiscal
outcomes. Spatial HAC standard errors would likely be larger than the reported
block-bootstrap SEs; this limitation biases toward finding significance and thus
works against the null result reported here.

\textbf{Control pool composition and smoke contamination.} By expanding the
never-treated group to include counties with low-to-moderate-severity fires,
this design compares counties along a common fire-exposure continuum. The primary
identification concern shifts from comparing fire-exposed to never-exposed counties
(the any-fire design) to comparing high-severity to low/moderate-severity fire
counties. A potential confound arises if low-severity fire counties also experienced
fiscal impacts from their fires, attenuating the estimated ATT. In this case, the
estimates represent a lower bound on the fiscal impact of high-severity relative to
low-severity fire---a locally identified but policy-relevant parameter, since many
risk-reduction interventions (prescribed burns, fuel treatments) aim to reduce
fire severity rather than fire occurrence. The effective sample size of the
reweighted control group is 19.6 counties out of 98, a constraint on statistical
power but an improvement over the any-fire specification (ESS = 11.1 of 51).

%% ============================================================
\section{Results}
%% ============================================================

\subsection{Extensive Margin: Fire Incidence}

Table~\ref{tab:main_ext} reports doubly-robust C\&S ATT estimates for the extensive
margin, comparing 293 counties that experienced a first qualifying wildfire
($\geq\!1{,}000$ acres) in 2013--2021 to 51 never-fired control counties (ESS = 11.1).
Column (1), the headline outcome, shows an ATT of \$\!-649 per capita on fiscal
balance (SE = 525, $p = 0.217$; 95\% CI: [\$\!-1{,}678, \$381]). The 2012
pre-treatment median fiscal balance for treated counties was approximately
\$1,580 per capita, so the point estimate represents a 41\% reduction---economically
large but statistically imprecise. No outcome achieves conventional significance.

Revenue declines are broad and large in magnitude. Total revenue per capita shows
ATT = \$\!-650 (SE = 646, $p = 0.314$). The largest point estimate is for
intergovernmental transfers: ATT = \$\!-841 (SE = 661, $p = 0.203$).
The negative sign is counterintuitive---post-fire counties are expected to
receive more federal and state aid, not less---but several mechanisms are
consistent with it. First, federal land payments (PILT and Secure Rural Schools),
which are tied to federal land area and not fire-triggered, trend separately
from emergency transfers; if treated (high-WFP) counties experienced declining
land payments relative to zero-fire controls over this period, the IGT ATT
would be negative. Second, the CoG IGT category aggregates state and federal
transfers across all programs, including formula-driven payments that may
grow more slowly in high-fire counties if their eligible populations (e.g.,
road miles, school enrollment) stagnate. Given the large standard error,
the estimate is statistically consistent with zero. Property tax revenue is near zero (\$+46, $p = 0.824$), consistent
with fire effects on assessed values being modest at this horizon across the
broad set of all fire-affected counties. Capital expenditures show
ATT = \$\!-226 (SE = 163, $p = 0.166$; 95\% CI: [\$\!-545, \$93])---the sharpest
signal, consistent with post-fire fiscal adjustment through deferred investment.
Total expenditures (\$\!-2, $p = 0.996$) and long-term debt (\$+55, $p = 0.869$)
show negligible effects.

The wide confidence intervals in the extensive-margin specification reflect the
small never-treated control group (51 counties, ESS = 11.1 after reweighting).
The covariate imbalance between fire-affected and never-affected counties is also
greater than in the intensive-margin design: pre-treatment WFP is 1,064 for
treated counties vs.\ 208 for controls (SMD = 0.89 unweighted), reflecting
systematic geographic sorting between high-hazard and fire-free counties.

\subsection{Intensive Margin: Fire Severity}

Table~\ref{tab:main_int} reports doubly-robust C\&S ATT estimates for the intensive
margin, comparing 246 counties that experienced a first high-severity wildfire
in 2013--2021 to 98 counties without a qualifying high-severity fire (ESS = 19.6).
Column (1) shows an ATT of \$\!-278 per capita on fiscal balance
(SE = 245, $p = 0.258$; 95\% CI: [\$\!-758, \$203]). The magnitude is
economically meaningful---an 18\% reduction from the 2012 median---but falls
short of conventional significance. No outcome achieves statistical significance.

Revenue falls broadly. Total revenue per capita shows ATT = \$\!-325
(SE = 306, $p = 0.289$). Intergovernmental transfers show ATT = \$\!-99
(SE = 148, $p = 0.504$), consistent with incomplete fiscal buffering through
federal and state aid. Property tax revenue is negative (\$\!-167, $p = 0.464$),
a stronger signal than in the extensive-margin design, consistent with
high-severity fire damage reducing market values enough to affect assessed values
over the four-to-nine-year observation horizon. Capital expenditures show
ATT = \$\!-181 (SE = 124, $p = 0.142$; 95\% CI: [\$\!-423, \$61])---the
sharpest estimate across outcomes, consistent with deferred capital investment
as the primary post-fire fiscal adjustment mechanism.

Comparing the two specifications, point estimates for fiscal balance are larger
in magnitude under the extensive margin (\$\!-649) than the intensive margin
(\$\!-278), but the extensive-margin SEs are also larger, reflecting the smaller
and less comparable control group. The intensive-margin design provides a cleaner
identification environment: pre-treatment parallel trends pass more strongly
(minimum $p = 0.50$ across outcomes vs.\ borderline results for some outcomes
in the extensive-margin design), and---as shown in Section~\ref{sec:robustness}---all
four estimators agree on negative signs for the intensive-margin specification
but not for the extensive-margin specification.

\subsection{Cohort Heterogeneity}
\label{sec:cohort}

Table~\ref{tab:cohort} disaggregates the ATT by treatment cohort for both margins.
Panel~A (extensive margin) shows that g\,=\,2017 counties carry a fiscal balance ATT
of \$\!-734 (SE = 607, $p = 0.227$) and capital outlays ATT of \$\!-283 (SE = 191,
$p = 0.139$). This cohort-level estimate averages two group-time cells:
ATT$(g{=}2017, t{=}2017) = {-}\$556$ (SE = 704, $p = 0.430$)---the clean
first-fire period---and ATT$(g{=}2017, t{=}2022) = {-}\$913$ (SE = 632, $p = 0.149$).
An important caveat applies to the latter cell: 192 of 208 g\,=\,2017 counties
(92.3\%) also experienced qualifying fires in 2017--2021, so ATT$(g{=}2017, t{=}2022)$
conflates the effect of the original 2013--2016 fire with cumulative subsequent
fire exposure. The clean first-fire estimate---ATT$(g{=}2017, t{=}2017)$---is
consistent with the headline null result but imprecise (SE = 704, wide interval
[\$\!-1{,}936, \$824]).
For g\,=\,2022 counties, fiscal balance is \$\!-211 (SE = 187,
$p = 0.258$) and capital outlays are near zero (\$+65, $p = 0.270$).
Panel~B (intensive margin) shows g\,=\,2017 fiscal balance ATT of \$\!-160
(SE = 220, $p = 0.467$) and capital outlays of \$\!-215 (SE = 143, $p = 0.132$).
The g\,=\,2022 cohort is observed in only one post-treatment census year (2022);
its ATT is therefore identified from a single cross-sectional comparison
between g\,=\,2022 counties and never-treated counties adjusted for covariates
rather than from within-unit pre-post variation. The resulting estimate for
fiscal balance (\$\!-865, SE = 1,420) is very imprecise and should be
interpreted with caution. Across both designs, the capital outlay
reduction is concentrated in g\,=\,2017 counties, which have two post-treatment
census years to absorb and reveal medium-run fiscal adjustment.

\subsection{Event Study}

Figure~\ref{fig:event_study} plots dynamic ATT estimates by event time
$\ell \in \{-10, -5, 0, 5\}$ for fiscal balance and capital expenditures under
the intensive-margin specification.
Pre-treatment estimates at $\ell = -10$ and $\ell = -5$ are statistically
indistinguishable from zero for both outcomes, consistent with pre-treatment
parallel trends. The post-treatment period ($\ell = 0$) shows negative point
estimates for both outcomes, with the capital expenditure decline materializing
at $\ell = 0$ and persisting at $\ell = 5$.
The joint pre-trend test is not rejected for any headline outcome
(all $p > 0.50$; range 0.50--0.99 across outcomes), providing strong
pre-treatment parallel trends evidence for the intensive-margin design.

\subsection{Dose-Response: Does Fire Size Predict Fiscal Effects?}
\label{sec:dose_response}

The preceding results compare counties that experienced any qualifying fire against
those that did not. A natural follow-on question is whether the null result holds
across the fire size distribution, or whether larger fires generate detectably worse
fiscal outcomes. To examine this, I construct a dose variable,
$D_i = \log(\text{MaxAcres}_i)$, equal to the log of the largest qualifying
MTBS fire experienced by county $i$ within its treatment cohort window
(2013--2016 for cohort $g = 2017$; 2017--2021 for cohort $g = 2022$).
The median $D_i$ among treated counties corresponds to approximately 13,161 acres.

The primary dose-response estimator is a two-way fixed effects regression,
\begin{equation}
  Y_{it} = \alpha_i + \gamma_t + \beta \, (D_i \times \text{Post}_{it}) + \varepsilon_{it},
  \label{eq:dose_twfe}
\end{equation}
where $\text{Post}_{it} = \mathbf{1}[\text{year\_cog}_t \geq g_i]$ for treated counties
and zero for never-treated counties. Unit and year fixed effects absorb all
time-invariant county characteristics and common time trends. Standard errors are
clustered by county. The coefficient $\beta$ captures the additional post-treatment
change in fiscal outcomes per unit increase in $\log(\text{MaxAcres})$, relative to
the never-treated baseline. This TWFE specification is subject to the same
heterogeneous-treatment-effects bias as standard TWFE in staggered designs
\citep{baker2022}: if dose-response slopes differ between the g\,=\,2017 and
g\,=\,2022 cohorts, the estimator may assign implicit negative weights to some
cohort-period combinations. Panel~B presents Callaway-Sant'Anna dose-bin estimates
as a complementary check that avoids this concern. A 10-fold increase in max fire size (e.g., 10,000 to
100,000 acres) implies an effect of $\log(10) \times \hat{\beta} \approx 2.3\,\hat{\beta}$.

Panel A of Table~\ref{tab:dose_response} presents TWFE estimates for all seven
fiscal outcomes. No coefficient is statistically significant. The fiscal balance
estimate is $\hat{\beta} = 10.76$ (SE = 11.56, $p = 0.353$), implying that a
10-fold increase in max fire size is associated with approximately \$25 per capita
in additional fiscal balance change---economically negligible against a pre-treatment
median of approximately \$4,200 per capita. Capital expenditures and total revenue
are similarly uninformative (both $p > 0.39$). Pre-trend Wald tests on the
year-by-dose interactions at $t = \{2002, 2007\}$ pass for all primary outcomes
(fiscal balance: $p = 0.357$; capital expenditures: $p = 0.188$). Property tax
revenue is the single exception (Wald $p = 0.019$), suggesting pre-existing
differences in property tax trajectories between high- and low-fire-risk counties
that pre-date the treatment window; the dose-response TWFE estimate for that
outcome should be interpreted with caution.

Panel B provides a complementary C\&S doubly-robust check. I split treated counties
at the median $D_i$ into low-dose ($n = 147$) and high-dose ($n = 146$) bins and
estimate the simple ATT for each bin against never-treated controls.
Low-dose counties generate a precisely null result for fiscal balance
(ATT = \$\!-179, SE = \$415, $p = 0.667$). High-dose counties produce a larger
negative estimate (ATT = \$\!-1{,}075, SE = \$635, $p = 0.091$) that is
marginally significant at the 10\% level but does not survive the conventional
5\% threshold. Both bins produce insignificant capital expenditure estimates.
The marginal high-dose result likely reflects concentration among large-fire
California counties; given the small treated-county effective sample and
the 5.4\% outlier rate in this sample, I do not interpret it as conclusive
evidence of a dose-response relationship.

Taken together, the dose-response analysis strengthens the main finding. The null
is not an artifact of pooling large and small fires: TWFE coefficients are
economically and statistically negligible across the fire size distribution,
and the C\&S dose-bin design finds no significant effects below the median
and only a borderline result above it.

%% ============================================================
\section{Robustness}
%% ============================================================

\subsection{Alternative Control Group and Sample Restrictions}

Table~\ref{tab:robustness} presents six robustness specifications for fiscal balance,
total revenue, and capital expenditures. The baseline specification (row 0) is
the main result. Row 1 expands the control group to include not-yet-treated
counties (g\,=\,2022 counties serve as controls for g\,=\,2017 through 2017),
shifting the fiscal balance estimate to \$\!-379 (SE = 448, $p = 0.397$).
The modest increase in magnitude with not-yet-treated controls---rather than
the attenuation expected under treatment effect contamination---suggests that
the g\,=\,2022 counties do not have materially different pre-treatment fiscal
trends from the never-treated pool.

Row 2 drops outlier counties (those with total revenue per capita more than
three times the state-year median) as complete observations to maintain a
balanced panel. Estimates are virtually unchanged from baseline.
Row 3 restricts the sample to six states (excluding Colorado and Utah) to
avoid the Utah NaN imputation for uninsurance and population density covariates;
estimates are stable at \$\!-278 (SE = 237, $p = 0.242$).
Row 4 substitutes the WHP 2014 raster for WFP 2012 in the propensity score;
the fiscal balance estimate shifts modestly to \$\!-263 (SE = 249, $p = 0.292$).
Row 5 excludes all California counties ($n = 52$ treated, $n = 5$ control counties
removed), testing whether findings reflect California's distinctive Proposition~13
property tax environment. Fiscal balance attenuates modestly to \$\!-209 (SE = 299,
$p = 0.485$); total revenue \textit{strengthens} to \$\!-422 (SE = 337, $p = 0.210$).
The strengthening of the revenue estimate upon California exclusion is consistent
with Proposition~13's reassessment-at-sale mechanism partially offsetting revenue
losses for California counties relative to non-California counties.
All six specifications are directionally negative and statistically insignificant.

The smoke exclusion buffer (100 km) was not applied as a sample restriction because
40 of 51 extensive-margin never-treated counties and 84 of 98 intensive-margin
never-treated counties are within 100 km of a fire perimeter in post-treatment
years. Restricting to the 11 (extensive) or 14 (intensive) smoke-free counties
would make estimation infeasible. For county fiscal outcomes, the primary
transmission channels---property assessment and debt issuance---are geographically
localised to the fire perimeter rather than the smoke plume; the relevant spillover
concern (smoke-driven productivity effects on local economic activity) operates
more directly through health and labor-market channels, which are not the focus
of this analysis.

\subsection{Alternative Estimators (Intensive Margin)}
\label{sec:robustness}

Table~\ref{tab:estimators} compares four estimators for the intensive-margin
specification. The C\&S doubly-robust estimate (column 1) of fiscal balance is
\$\!-278. Sun and Abraham \citeyearpar{sun2021} yield \$\!-126 (SE = 187,
$p = 0.499$). Synthetic DiD \citep{arkhangelsky2021} yields \$\!-252 (SE = 202,
$p = 0.211$). Two-stage DiD \citep{gardner2022} yields \$\!-188 (SE = 222,
$p = 0.397$).

All four estimators agree on a negative sign for fiscal balance and for total
and property tax revenue. This cross-estimator consensus is attributable to the
closer counterfactual provided by the expanded low-severity-fire control group,
which reduces the observable selection problem that drives sign disagreement
in the extensive-margin design. In the extensive-margin specification
(any qualifying fire vs.\ zero-fire counties), the larger covariate imbalance
causes non-DR estimators to yield positive estimates for fiscal balance, while
the DR estimator---which adjusts for observable selection---yields negative
estimates. The intensive-margin design, with a more comparable control group,
produces consistent negative estimates across all estimators.

Capital expenditures diverge: the C\&S DR estimate (\$\!-181) is negative while
the three non-DR estimators are near-zero positive (\$6 to \$23). The DR
adjustment accounts for higher pre-treatment capital spending baselines in
treated counties relative to the expanded control group.

%% ============================================================
\section{Discussion and Conclusion}
%% ============================================================

This paper provides causal estimates of wildfire fiscal effects on county
governments across eight western states, extending prior California evidence to
county governments with diverse property tax institutions. Using Census of
Governments quinquennial data from 2002 to 2022 and WFP-matched staggered DiD
designs, I find uniformly negative but statistically insignificant effects on
fiscal balance across both the extensive margin (fire incidence) and the intensive
margin (fire severity). Capital expenditures provide the most consistent signal (\$\!-226 per
capita, extensive; \$\!-181 per capita, intensive), consistent with deferred capital
investment as the primary post-fire fiscal adjustment channel.

Three substantive conclusions emerge. First, neither the incidence nor the severity
of wildfires generates a detectable fiscal effect at conventional significance levels
within the quinquennial observation window. The null result is informative: point
estimates suggest economically meaningful deterioration (18--41\% of the pre-treatment
median fiscal balance), and the wide confidence intervals are consistent with this
range. Second, intergovernmental transfers do not appear to offset local revenue
losses---the IGT point estimates are negative under both designs---mirroring the
finding that the federal wildfire aid architecture provides less automatic fiscal
buffering than hurricane-specific programs. Third, the intensive-margin design
demonstrates a methodological advantage over the extensive-margin design: the
expanded low-to-moderate-severity control group is a more credible counterfactual,
reduces observable selection, and yields cross-estimator sign agreement that validates
the C\&S doubly-robust estimate as internally consistent.

A dose-response analysis confirms that the null is not an artifact of pooling
large and small fires. TWFE estimates of the effect of log maximum fire size on
all seven fiscal outcomes are statistically and economically negligible
(fiscal balance: \$11 per log-acre, SE = \$12, $p = 0.35$; IQR-implied effect
$\approx$\,\$27 per capita). Callaway-Sant'Anna dose-bin estimates for counties
below and above the median fire size (13,161 acres) are both insignificant for
capital expenditures; a borderline negative estimate for high-dose fiscal balance
($p = 0.091$) does not survive the conventional threshold and is sensitive to
California outlier counties.

The primary limitation is statistical power. The effective sample sizes after
doubly-robust reweighting are 11.1 counties (extensive) and 19.6 counties (intensive);
confidence intervals are consistent with effects ranging from zero to losses exceeding
major hurricane estimates. Table~\ref{tab:mde} quantifies the inferential limits:
the minimum detectable effect at 80\% power ($\alpha = 0.05$) is \$687 per capita
for the intensive-margin fiscal balance outcome and \$1{,}471 for the extensive
margin---seven and fifteen times the \citet{liao2022} California benchmark of
approximately \$97 per capita, respectively. For capital outlays, the MDEs are
\$346 (intensive) and \$456 (extensive). These bounds clarify that the null
result reflects insufficient power rather than confirmed fiscal resilience:
effects of the magnitude found in the California literature remain within the
confidence intervals for every outcome. Annual CoG survey data---with more
post-treatment observation windows---would substantially improve precision. The negative
property-tax revenue estimate in the intensive-margin design (\$\!-167 per capita)
contrasts with the near-zero extensive-margin result, consistent with high-severity
fire damage reducing market values enough to affect assessed values at the four-to-nine
year horizon. This contrasts with the \citet{liao2022} California result, where
Proposition 13's reassessment-at-sale mechanism produces revenue increases following
fire-induced sales. Understanding which fiscal institutions---property tax assessment
rules, state aid formulas, and home rule authority---mediate the wildfire fiscal shock
across states is the natural next step.

\clearpage

%% ============================================================
%% REFERENCES
%% ============================================================

\bibliographystyle{chicago}
\begin{thebibliography}{99}

\bibitem[Arkhangelsky et al.(2021)]{arkhangelsky2021}
Arkhangelsky, D., S.~Athey, D.~A.~Hirshberg, G.~W.~Imbens, and S.~Wager.
\newblock 2021.
\newblock Synthetic difference-in-differences.
\newblock \textit{American Economic Review} 111(12): 4088--4118.

\bibitem[Baker et al.(2022)]{baker2022}
Baker, A.~C., D.~F.~Larcker, and C.~C.~Y.~Wang.
\newblock 2022.
\newblock How much should we trust staggered difference-in-differences estimates?
\newblock \textit{Journal of Financial Economics} 144(2): 370--395.

\bibitem[Baylis and Boomhower(2023)]{baylis2023}
Baylis, P. and J.~Boomhower.
\newblock 2023.
\newblock The economic incidence of wildfire suppression in the United States.
\newblock \textit{American Economic Journal: Applied Economics} 15(1): 442--473.

\bibitem[Borgschulte, Molitor, and Zou(2024)]{borgschulte2024}
Borgschulte, M., D.~Molitor, and E.~Y.~Zou.
\newblock 2024.
\newblock Air pollution and the labor market: Evidence from wildfire smoke.
\newblock \textit{Review of Economics and Statistics} 106(6): 1691--1708.

\bibitem[Callaway and Sant'Anna(2021)]{callaway2021}
Callaway, B. and P.~H.~C.~Sant'Anna.
\newblock 2021.
\newblock Difference-in-differences with multiple time periods.
\newblock \textit{Journal of Econometrics} 225(2): 200--230.

\bibitem[Chen and Cohen(2025)]{chen2025}
Chen, H. and L.~Cohen.
\newblock 2025.
\newblock Assessing assessors.
\newblock NBER Working Paper No.~33238.

\bibitem[Deryugina(2017)]{deryugina2017}
Deryugina, T.
\newblock 2017.
\newblock The fiscal cost of hurricanes: Disaster aid versus social insurance.
\newblock \textit{American Economic Journal: Economic Policy} 9(3): 168--198.

\bibitem[Gardner(2022)]{gardner2022}
Gardner, J.
\newblock 2022.
\newblock Two-stage differences in differences.
\newblock Working paper.

\bibitem[Jerch, Kahn, and Lin(2023)]{jerch2023}
Jerch, R., M.~E.~Kahn, and G.~C.~Lin.
\newblock 2023.
\newblock Local public finance dynamics and hurricane shocks.
\newblock \textit{Journal of Urban Economics} 134: 103521.

\bibitem[Liao and Kousky(2022)]{liao2022}
Liao, Y. and C.~Kousky.
\newblock 2022.
\newblock The fiscal impacts of wildfires on California municipalities.
\newblock \textit{Journal of the Association of Environmental and Resource Economists}
  9(3): 455--493.

\bibitem[Lutz, Molloy, and Shan(2011)]{lutz2011}
Lutz, B., R.~Molloy, and H.~Shan.
\newblock 2011.
\newblock The housing crisis and state and local government tax revenue: Five channels.
\newblock \textit{Regional Science and Urban Economics} 41(4): 306--319.

\bibitem[Roth(2022)]{roth2022}
Roth, J.
\newblock 2022.
\newblock Pre-test with caution: Event-study estimates after testing for parallel trends.
\newblock \textit{American Economic Review: Insights} 4(3): 305--322.

\bibitem[Scott et al.(2013)]{scott2013}
Scott, J.~H., M.~P.~Thompson, and D.~E.~Calkin.
\newblock 2013.
\newblock A wildfire risk assessment framework for land and resource management.
\newblock USDA Forest Service General Technical Report RMRS-GTR-315.

\bibitem[Sun and Abraham(2021)]{sun2021}
Sun, L. and S.~Abraham.
\newblock 2021.
\newblock Estimating dynamic treatment effects in event studies with heterogeneous
  treatment effects.
\newblock \textit{Journal of Econometrics} 225(2): 175--199.

\end{thebibliography}

\clearpage

%% ============================================================
%% TABLES
%% ============================================================

\begin{table}[h!]
\begin{threeparttable}
\caption{Covariate Balance: Treated (High-Severity Fire) vs.\ Never-Treated Counties (2012)}
\label{tab:balance}
\centering
\footnotesize
\begin{tabular}{lrrrrrrr}
\toprule
 & \multicolumn{3}{c}{Unweighted} & \phantom{abc} &
   \multicolumn{3}{c}{IPW-Weighted} \\
\cmidrule{2-4} \cmidrule{6-8}
Covariate & Treated & Control & SMD & & Treated & Control & SMD \\
\midrule
\textit{Wildfire hazard} \\
\quad WFP 2012 (mean)        & 1064.0 & 207.6 & 0.89 & & 1064.0 &  903.4 & 0.13 \\
\quad WFP quintile           &   3.46 &   1.91 & 1.32 & &   3.46 &    3.28 & 0.14 \\
\quad Pre-2013 fire (=1)     &  0.923 &  0.561 & 0.91 & &  0.923 &   0.929 &$-$0.03 \\
\quad Pre-2013 log acres     & 49.5  &  15.6 & 1.20 & & 49.5  &   43.4 & 0.19 \\
\addlinespace
\textit{Geography \& demographics} \\
\quad RUCC 2013              &   5.29 &   5.51 &$-$0.08& &   5.29 &    5.18 & 0.04 \\
\quad Log median HH income   & 10.90 & 10.90 & 0.04& &10.90 &  10.90 & 0.04\\
\quad Poverty rate (\%)      & 14.4  &  14.1 & 0.06 & & 14.4  &   14.4 & 0.00 \\
\quad Uninsurance rate (\%)  & 17.3  &  16.7 & 0.13 & & 17.3  &   18.1 &$-$0.15\\
\quad Log pop.\ density      &  2.06 &  2.24 &$-$0.11& & 2.06 &   2.38 &$-$0.22\\
\addlinespace
\textit{Pre-treatment fiscal (2012)} \\
\quad Log revenue p.c.       &  7.42 &  7.39 & 0.03 & &  7.42 &   7.37 & 0.07 \\
\quad Log property tax p.c.  &  6.11 &  6.10 & 0.02 & &  6.11 &   6.00 & 0.16 \\
\quad Log LT debt p.c.       &  4.14 &  4.22 &$-$0.03& & 4.14 &   4.31 &$-$0.06\\
\midrule
Observations                 & \multicolumn{3}{c}{$N_{\rm trt}=246$, $N_{\rm ctl}=98$} & &
                               \multicolumn{3}{c}{ESS = 19.6 of 98} \\
\bottomrule
\end{tabular}
\begin{tablenotes}[flushleft]
\small
\item \textit{Notes}: Treated = counties whose first qualifying MTBS wildfire
in 2013--2021 had a detectable High-severity dNBR class (\texttt{HIGH\_THRES}
between 100 and 5,000 dNBR units). Control = counties without a qualifying
high-severity fire in 2013--2021, including counties that experienced only
low-to-moderate-severity fires. Standardized mean differences (SMD) computed as
(treated mean $-$ control mean) / pooled standard deviation.
IPW weights: treated counties $w=1$; control counties $w=\hat{p}/(1-\hat{p})$,
normalized so control weights sum to the number of control units.
Effective sample size (ESS) for reweighted controls = $[\sum w_i]^2 / \sum w_i^2$.
Uninsurance rate and log population density imputed to state-year median
for 29 Utah counties absent from ACS 5-year estimates prior to 2009;
for groups where state-year median is missing, the overall sample median is used.
\end{tablenotes}
\end{threeparttable}
\end{table}

\clearpage

%% ── TABLE 2: EXTENSIVE MARGIN ────────────────────────────────────────────────

\begin{table}[h!]
\begin{threeparttable}
\caption{Main Results: Doubly-Robust C\&S ATT Estimates --- Extensive Margin (Any Qualifying Fire)}
\label{tab:main_ext}
\centering
\small
\begin{tabular}{lrrrrrrr}
\toprule
 & (1) & (2) & (3) & (4) & (5) & (6) & (7) \\
 & Fiscal & Total & Total & Property & Capital & IGT & LT  \\
 & balance & revenue & expend. & tax rev. & outlays & transfers & debt \\
\midrule
ATT (\$2019 per capita)
  & $-$648.5 & $-$650.1 & $-$1.6   & $+$46.1   & $-$225.9 & $-$841.0 & $+$54.5 \\
SE
  & (525.2) & (645.9) & (345.0) & (207.5) & (162.9)  & (660.6) & (330.0)  \\
95\% CI
  & [$-$1678, & [$-$1916, & [$-$678, & [$-$361, & [$-$545, & [$-$2136, & [$-$592, \\
  & \phantom{[}381] & \phantom{[}616] & \phantom{[}675] & \phantom{[}453] & \phantom{[}93] & \phantom{[}454] & \phantom{[}701] \\
$p$-value
  & 0.217   & 0.314   & 0.996   & 0.824   & 0.166    & 0.203   & 0.869    \\
\addlinespace
Treated counties & 293 & 293 & 293 & 293 & 293 & 293 & 293 \\
Control counties &  51 &  51 &  51 &  51 &  51 &  51 &  51 \\
Observations & 1,720 & 1,720 & 1,720 & 1,720 & 1,720 & 1,720 & 1,720 \\
\bottomrule
\end{tabular}
\begin{tablenotes}[flushleft]
\small
\item \textit{Notes}: Callaway and Sant'Anna (2021) doubly-robust group-time ATT
estimator. Treatment: county's first qualifying MTBS wildfire ($\geq\!1{,}000$ acres)
in 2013--2021, regardless of fire severity. Control group: counties with no
qualifying MTBS wildfire in 2013--2021 (never-fired; $n = 51$).
Covariates, estimation method, and SE computation identical to Table~\ref{tab:main_int}.
ESS of reweighted controls = 11.1.
IGT = intergovernmental transfers; LT debt = long-term debt.
All outcomes in 2019\$ per capita. *** $p < 0.01$; ** $p < 0.05$; * $p < 0.10$.
\end{tablenotes}
\end{threeparttable}
\end{table}

\clearpage

%% ── TABLE 3: INTENSIVE MARGIN ────────────────────────────────────────────────

\begin{table}[h!]
\begin{threeparttable}
\caption{Main Results: Doubly-Robust C\&S ATT Estimates --- Intensive Margin (High-Severity Fire)}
\label{tab:main_int}
\centering
\small
\begin{tabular}{lrrrrrrr}
\toprule
 & (1) & (2) & (3) & (4) & (5) & (6) & (7) \\
 & Fiscal & Total & Total & Property & Capital & IGT & LT  \\
 & balance & revenue & expend. & tax rev. & outlays & transfers & debt \\
\midrule
ATT (\$2019 per capita)
  & $-$277.5 & $-$324.7 & $-$47.3   & $-$166.9   & $-$181.4 & $-$99.1 & 727.5 \\
SE
  & (245.3) & (306.3) & (294.5) & (227.8) & (123.5)  & (148.2) & (960.0)  \\
95\% CI
  & [$-$758, & [$-$925, & [$-$625, & [$-$613, & [$-$423, & [$-$390, & [$-$1154, \\
  & \phantom{[}203] & \phantom{[}276] & \phantom{[}530] & \phantom{[}280] & \phantom{[}61] & \phantom{[}191] & \phantom{[}2609] \\
$p$-value
  & 0.258   & 0.289   & 0.873   & 0.464   & 0.142    & 0.504   & 0.449    \\
\addlinespace
Treated counties & 246 & 246 & 246 & 246 & 246 & 246 & 246 \\
Control counties &  98 &  98 &  98 &  98 &  98 &  98 &  98 \\
Observations & 1,720 & 1,720 & 1,720 & 1,720 & 1,720 & 1,720 & 1,720 \\
\bottomrule
\end{tabular}
\begin{tablenotes}[flushleft]
\small
\item \textit{Notes}: Callaway and Sant'Anna (2021) doubly-robust group-time
ATT estimator. Treatment: counties whose first qualifying MTBS wildfire in
2013--2021 had a detectable High-severity dNBR class (\texttt{HIGH\_THRES}
between 100 and 5,000 dNBR units). Control group: counties without a qualifying
high-severity fire in 2013--2021, including counties with only low-to-moderate
severity fires. Covariates in propensity score and outcome regression: WFP 2012
quintile, pre-2013 fire indicator, log pre-2013 cumulative fire acres, RUCC 2013,
log median household income, poverty rate, uninsurance rate (imputed for UT),
log population density (imputed for UT). Estimation: \texttt{did} package v2.3.0,
\texttt{base\_period = ``varying''}, doubly-robust. Standard errors in parentheses;
95\% confidence intervals in brackets. Both computed via block bootstrap with 999
replications, clustered by county. ESS of reweighted controls = 19.6.
IGT = intergovernmental transfers (federal + state); LT debt = long-term debt.
All outcomes in 2019\$ per capita, deflated using national CPI-U.
*** $p < 0.01$; ** $p < 0.05$; * $p < 0.10$.
\end{tablenotes}
\end{threeparttable}
\end{table}

\clearpage

%% ── TABLE 4: COHORT ATTs FOR BOTH MARGINS ────────────────────────────────────

\begin{table}[h!]
\begin{threeparttable}
\caption{Cohort-Specific ATTs: Both Margins (g\,=\,2017 vs.\ g\,=\,2022)}
\label{tab:cohort}
\centering
\small
\begin{tabular}{lrrrrrr}
\toprule
 & (1) & (2) & (3) & (4) & (5) & (6) \\
 & Fiscal & Total & Total & Property & Capital & IGT \\
 & balance & revenue & expend. & tax rev. & outlays & transfers \\
\midrule
\textit{\underline{Panel A: Extensive margin} (any fire $\geq$1,000 ac)} \\
\addlinespace
\textit{\quad Early cohort (g\,=\,2017; first fire 2013--2016; $n=208$)} \\
ATT & $-$734.4 & $-$784.4 & ---  & ---  & $-$282.9 & --- \\
SE  & (607.4)  & (765.9)  & ---  & ---  & (191.1)  & --- \\
$p$ & 0.227    & 0.306    & ---  & ---  & 0.139    & --- \\
\addlinespace
\textit{\quad Late cohort (g\,=\,2022; first fire 2017--2021; $n=85$)} \\
ATT & $-$211.0 & $+$34.0  & ---  & ---  & $+$64.6  & --- \\
SE  & (186.6)  & (274.6)  & ---  & ---  & (58.5)   & --- \\
$p$ & 0.258    & 0.901    & ---  & ---  & 0.270    & --- \\
\addlinespace\midrule\addlinespace
\textit{\underline{Panel B: Intensive margin} (high-severity fire)} \\
\addlinespace
\textit{\quad Early cohort (g\,=\,2017; sev.\ fires 2013--2016; $n=173$)} \\
ATT & $-$160.4 & $-$248.6 & $-$88.3  & $-$29.5  & $-$215.1 & $-$111.1 \\
SE  & (220.0)  & (349.0)  & (383.4)  & (113.1)  & (142.6)  & (174.4)  \\
$p$ & 0.467    & 0.476    & 0.818    & 0.794    & 0.132    & 0.524    \\
\addlinespace
\textit{\quad Late cohort (g\,=\,2022; sev.\ fires 2017--2021; $n=73$)} \\
ATT & $-$864.9 & $-$706.5  & 158.4   & $-$856.1  & $-$12.3     & $-$39.3 \\
SE  & (1419.9)  & (1355.3) & (281.8) & (1339.1) & (43.2)   & (57.1)  \\
$p$ & 0.542    & 0.602   & 0.574   & 0.522  & 0.776    & 0.492   \\
\midrule
Control (extensive) & 51 & 51 & 51 & 51 & 51 & 51 \\
Control (intensive) & 98 & 98 & 98 & 98 & 98 & 98 \\
\bottomrule
\end{tabular}
\begin{tablenotes}[flushleft]
\small
\item \textit{Notes}: Group-level ATT aggregations from the C\&S doubly-robust
estimator. Panel A: extensive-margin cohorts (any qualifying fire). Panel B:
intensive-margin cohorts (high-severity fire). Within each panel, early-cohort
estimates average ATT$(g{=}2017, t{=}2017)$ and ATT$(g{=}2017, t{=}2022)$ with
cohort-size weights; counties observed in two post-treatment census years.
Late-cohort estimates equal ATT$(g{=}2022, t{=}2022)$; observed in one post-treatment
year and identified from a single cross-sectional comparison adjusted for covariates.
Dashes indicate outcomes not estimated in the cohort decomposition run.
\textit{Double-treatment caveat (extensive margin, g\,=\,2017):} 192 of 208
early-cohort counties (92.3\%) also experienced qualifying fires in 2017--2021.
The g\,=\,2017 cohort ATT therefore averages ATT$(g{=}2017, t{=}2017) = {-}\$556$
(SE = 704; clean first-fire effect) with ATT$(g{=}2017, t{=}2022) = {-}\$913$
(SE = 632; confounded by subsequent fire exposure). See Section~\ref{sec:cohort}.
Sev.\ = high-severity.
\end{tablenotes}
\end{threeparttable}
\end{table}

\clearpage

\begin{table}[h!]
\begin{threeparttable}
\caption{Robustness Checks: Doubly-Robust C\&S ATT across Specifications}
\label{tab:robustness}
\centering
\small
\begin{tabular}{lrrrrrr}
\toprule
 & \multicolumn{2}{c}{(1) Fiscal balance} &
   \multicolumn{2}{c}{(2) Total revenue} &
   \multicolumn{2}{c}{(3) Capital outlays} \\
\cmidrule(lr){2-3}\cmidrule(lr){4-5}\cmidrule(lr){6-7}
Specification & ATT & (SE) & ATT & (SE) & ATT & (SE) \\
\midrule
(0) Baseline              & $-$277.5 & (246.2) & $-$324.7 & (318.7) & $-$181.4 & (115.6) \\
(1) Not-yet-treated ctrl  & $-$379.4 & (448.1) & $-$361.7 & (460.5) & $-$118.5 & (78.0)  \\
(2) Drop outlier counties & $-$277.5 & (251.4) & $-$324.7 & (321.2) & $-$181.4 & (121.8) \\
(3) 6-state (excl.\ CO, UT) & $-$277.5 & (236.9) & $-$324.7 & (321.9) & $-$181.4 & (114.2)\\
(4) WHP 2014 in $\mathbf{X}$   & $-$262.9 & (249.4) & $-$335.5 & (345.3) & $-$154.9 & (106.8) \\
(5) Excl.\ California     & $-$208.7 & (299.0) & $-$422.3 & (336.8) & $-$93.2  & (82.7)  \\
\midrule
Treated ($N$) & 246 & & 246 & & 246 & \\
Control ($N$) &  98 & &  98 & &  98 & \\
\footnotesize{(Row 5: $N_{\rm trt}=194$, $N_{\rm ctl}=93$)} & & & & & & \\
\bottomrule
\end{tabular}
\begin{tablenotes}[flushleft]
\small
\item \textit{Notes}: All specifications use doubly-robust C\&S estimator with
never-treated controls (except row 1) and bootstrap SE (999 replications).
Treatment defined as in Table~\ref{tab:main_int}.
(0) Baseline: same as Table~\ref{tab:main_int} column (1). (1) Not-yet-treated
controls: control group expands to include g\,=\,2022 counties as controls for
the g\,=\,2017 cohort in 2017. (2) Outlier counties---those with total revenue
per capita exceeding three times the state-year median in any year---are dropped
as complete observations to preserve panel balance.
(3) Drops Colorado and Utah; treated $n = 221$, control $n = 94$.
(4) WHP 2014 (USFS Wildfire Hazard Potential 2014) replaces WFP 2012 in the
covariate vector $\mathbf{X}$; note WHP 2014 is not predetermined for 2013--2014
fires. (5) Excludes all California counties ($n{=}52$ treated, $n{=}5$ control
removed), reducing the sample to 194 treated and 93 control counties; tests
whether results reflect California's Proposition~13 property tax environment.
None of the specifications yield statistical significance at conventional levels;
all six are directionally negative for fiscal balance and revenue. Under California
exclusion, total revenue strengthens to \$\!-422 (vs.\ \$\!-325 baseline),
consistent with Proposition~13's reassessment-at-sale mechanism partially
offsetting revenue losses for California counties.
\end{tablenotes}
\end{threeparttable}
\end{table}

\clearpage

\begin{table}[h!]
\begin{threeparttable}
\caption{Comparison of Heterogeneity-Robust Estimators (High-Severity Fire Treatment)}
\label{tab:estimators}
\centering
\small
\begin{tabular}{lrrrr}
\toprule
 & (1) & (2) & (3) & (4) \\
 & C\&S DR & Sun-Abraham & Synthetic DiD & Two-Stage DiD \\
\midrule
\textit{Panel A: Fiscal balance per capita} \\
ATT         & $-$277.5 & $-$126.4   & $-$252.2    & $-$188.3   \\
SE          & (245.3)  & (187.0) & (201.6) & (222.3) \\
95\% CI     & [$-$758, 203] & [$-$493, 240] & [$-$647, 143] & [$-$624, 248] \\
$p$-value   & 0.258    & 0.499   & 0.211   & 0.397   \\
\addlinespace
\textit{Panel B: Total revenue per capita} \\
ATT         & $-$324.7 & $-$15.7    & $-$137.6    & $-$161.8    \\
SE          & (306.3)  & (209.7) & (225.6) & (238.5) \\
$p$-value   & 0.289    & 0.940   & 0.542   & 0.498   \\
\addlinespace
\textit{Panel C: Property tax revenue per capita} \\
ATT         & $-$166.9     & $-$140.6    & $-$231.4    & $-$243.5    \\
SE          & (227.8)  & (175.9) & (177.6) & (215.7) \\
$p$-value   & 0.464    & 0.424   & 0.193   & 0.259   \\
\addlinespace
\textit{Panel D: Capital expenditures per capita} \\
ATT         & $-$181.4 & 6.2  & 16.1    & 23.1    \\
SE          & (123.5)  & (41.1)  & (20.4)  & (22.9)  \\
$p$-value   & 0.142    & 0.879   & 0.431   & 0.313   \\
\midrule
Covariate adjustment & DR & None & Synthetic & FE \\
Cohort coding & g\,=\,0/2017/2022 & g\,=\,0/2017/2022 & g\,=\,0/2017/2022 & g\,=\,0/2017/2022 \\
\bottomrule
\end{tabular}
\begin{tablenotes}[flushleft]
\small
\item \textit{Notes}: Treatment defined as in Table~\ref{tab:main_int} (high-severity fires only).
Column (1): Callaway-Sant'Anna (2021) doubly-robust ATT, block bootstrap SE (999 replications).
Column (2): Sun and Abraham (2021) heterogeneity-robust TWFE, \texttt{sunab(g, year)} in
\texttt{fixest}; cohorts g\,=\,0 (never-treated), 2017, 2022; base period $\ell = -5$;
average over $\ell \geq 0$ cells; clustered SE.
Column (3): Synthetic DiD \citep{arkhangelsky2021}; separate estimates for g\,=\,2017
(donors: g\,=\,0 + g\,=\,2022 not-yet-treated; $T_0=3$) and g\,=\,2022 (donors: g\,=\,0;
$T_0=4$); pooled with cohort-size weights; bootstrap SE (500 replications).
Column (4): Two-stage DiD \citep{gardner2022} via \texttt{did2s} package v1.0.2;
stage-1 absorbs county and year fixed effects from pre-treatment observations;
stage-2 regresses residuals on treatment indicator; clustered SE.
Unlike the all-fire specification where columns (2)--(4) yielded positive estimates,
all four estimators agree on negative signs for Panels A--C when treatment is
restricted to high-severity fires, reflecting improved counterfactual comparability.
Capital expenditures (Panel D) diverge: the DR estimate is negative ($-$\$181) while
non-DR estimators are near-zero positive, suggesting the DR covariate adjustment
accounts for higher pre-treatment capital spending baselines in treated counties.
*** $p < 0.01$; ** $p < 0.05$; * $p < 0.10$.
\end{tablenotes}
\end{threeparttable}
\end{table}

\clearpage

\begin{table}[h!]
\begin{threeparttable}
\caption{Dose-Response DiD --- Effect of Maximum Fire Size on Fiscal Outcomes}
\label{tab:dose_response}
\centering
\small

\textit{Panel A: TWFE dose-response (all 344 counties; treated $n=293$, never-treated $n=51$)} \\[4pt]
\begin{tabular}{lrrrr}
\toprule
Outcome & $\hat{\beta}$ & SE & $p$-value & Pre-trend Wald $p$ \\
\midrule
Fiscal balance p.c.          &  10.76 & (11.56) & 0.353 & 0.357 \\
Total revenue p.c.           &   3.95 & (13.23) & 0.765 & 0.565 \\
Property tax revenue p.c.    &   6.02 & (10.26) & 0.558 & 0.019\textsuperscript{$\dagger$} \\
Intergovernmental revenue p.c. & $-$0.72 & (4.56) & 0.875 & 0.158 \\
Total expenditures p.c.      &  $-$6.80 & (7.99) & 0.395 & 0.935 \\
Capital expenditures p.c.    &   1.56 &  (2.05) & 0.446 & 0.188 \\
Long-term debt p.c.          &   2.55 &  (9.11) & 0.780 & 0.970 \\
\bottomrule
\end{tabular}

\medskip
\textit{Panel B: C\&S DR dose-bin ATT (median split at $\log D_i = 9.49$, i.e.\ 13{,}161 acres)} \\[4pt]
\begin{tabular}{lrrrrrr}
\toprule
 & \multicolumn{3}{c}{Low dose ($\leq$ median, $n=147$)} & \multicolumn{3}{c}{High dose ($>$ median, $n=146$)} \\
\cmidrule(lr){2-4} \cmidrule(lr){5-7}
Outcome & ATT & SE & $p$ & ATT & SE & $p$ \\
\midrule
Fiscal balance p.c.       & $-$178.6 & (414.6) & 0.667 & $-$1{,}074.9 & (635.1) & 0.091\textsuperscript{*} \\
Capital expenditures p.c. & $-$118.3 &  (85.5) & 0.166 &   $-$260.1 & (201.3) & 0.196 \\
\bottomrule
\end{tabular}

\begin{tablenotes}[flushleft]
\small
\item \textit{Notes}: Dose variable $D_i = \log(\text{MaxAcres}_i)$ is the log of the largest
qualifying MTBS fire ($\geq 1{,}000$ acres) within the county's treatment cohort window
(2013--2016 for cohort $g{=}2017$; 2017--2021 for cohort $g{=}2022$).
Panel A: two-way fixed effects regression of equation~\eqref{eq:dose_twfe};
$\hat{\beta}$ is the coefficient on $D_i \times \text{Post}_{it}$; county and year fixed
effects included; SE clustered by county.
Pre-trend Wald $p$: joint Wald test on dose interactions at $t \in \{2002,2007\}$
(2 d.f.); clustered.
Panel B: Callaway-Sant'Anna (2021) doubly-robust simple ATT estimated separately
for counties below and above the median dose; never-treated controls; same
covariates as Table~\ref{tab:main_ext}.
Fiscal outcomes in real 2019 dollars per capita.
\textsuperscript{$\dagger$} Pre-trend test rejected at 5\%: property tax revenue
trajectories differ pre-treatment across fire-size groups; dose-response
coefficient for this outcome is not causally interpretable.
*** $p<0.01$; ** $p<0.05$; * $p<0.10$.
\end{tablenotes}
\end{threeparttable}
\end{table}

\clearpage

%% ============================================================
%% FIGURES (placeholders for actual PDF figures)
%% ============================================================

\begin{figure}[h!]
  \centering
  \caption{Event Study: Fiscal Balance and Capital Expenditures per Capita}
  \label{fig:event_study}
  \includegraphics[width=0.48\textwidth]{../output/figures/event_study_fiscal_balance_pc.pdf}
  \hfill
  \includegraphics[width=0.48\textwidth]{../output/figures/event_study_exp_capital_pc.pdf}
  \begin{minipage}{\textwidth}
  \footnotesize
  \textit{Notes}: Doubly-robust C\&S dynamic ATT estimates by event time $\ell$,
  measured in calendar years relative to the treatment cohort year $g$
  (quinquennial CoG spacing: $\ell \in \{-10, -5, 0, +5\}$ corresponds to
  census years $g-10, g-5, g, g+5$; for the primary g\,=\,2017 cohort,
  $\ell = -10, -5, 0, +5$ maps to 2007, 2012, 2017, 2022). Left panel:
  fiscal balance per capita. Right panel: capital expenditures per capita.
  Vertical lines are 95\% bootstrap confidence intervals (clustered by county).
  Dashed vertical line marks the pre/post divide. Pre-treatment joint
  $\chi^2$ test: $p > 0.05$ for both outcomes.
  \end{minipage}
\end{figure}

\clearpage

%% ── TABLE A1: MINIMUM DETECTABLE EFFECTS ─────────────────────────────────────

\begin{table}[h!]
\begin{threeparttable}
\caption{Statistical Power: Minimum Detectable Effects at 80\% Power ($\alpha = 0.05$)}
\label{tab:mde}
\centering
\small
\begin{tabular}{lrrrrrr}
\toprule
 & \multicolumn{3}{c}{Intensive margin (ESS\,=\,19.6)} &
   \multicolumn{3}{c}{Extensive margin (ESS\,=\,11.1)} \\
\cmidrule(lr){2-4}\cmidrule(lr){5-7}
Outcome & ATT & SE & MDE & ATT & SE & MDE \\
\midrule
Fiscal balance p.c.       & $-$278 & 245 & 687   & $-$649 & 525 & 1{,}471 \\
Total revenue p.c.        & $-$325 & 306 & 858   & $-$650 & 646 & 1{,}810 \\
Capital outlays p.c.      & $-$181 & 124 & 346   & $-$226 & 163 &   456   \\
Property tax revenue p.c. & $-$167 & 228 & 638   &  $+$46 & 208 &   581   \\
IGT transfers p.c.        &  $-$99 & 148 & 415   & $-$841 & 661 & 1{,}851 \\
Total expenditures p.c.   &  $-$47 & 295 & 825   &   $-$2 & 345 &   967   \\
LT debt p.c.              & $+$728 & 960 & 2{,}690 &  $+$55 & 330 &   925   \\
\addlinespace
\textit{Reference benchmark} \\
\quad Liao-Kousky (2022) CA fiscal balance & \multicolumn{6}{c}{$\approx$\,\$97 per capita} \\
\bottomrule
\end{tabular}
\begin{tablenotes}[flushleft]
\small
\item \textit{Notes}: MDE $= (z_{0.025} + z_{0.20}) \times \hat{\sigma} = 2.802 \times \hat{\sigma}$,
where $\hat{\sigma}$ is the block bootstrap SE from the C\&S doubly-robust estimator
(Tables~\ref{tab:main_ext} and \ref{tab:main_int}).
All values in 2019 USD per capita. ESS = effective sample size of reweighted
control group after doubly-robust reweighting. The MDE is the minimum true effect
that the design detects with 80\% probability at a 5\% significance level (two-sided).
Effects smaller than the MDE cannot be ruled out. The \citet{liao2022} California
municipal benchmark of approximately \$97 per capita falls below the MDE for every
outcome in both margins: the design would need approximately 50$\times$ (intensive)
to 230$\times$ (extensive) reductions in SE to reliably detect effects of that magnitude.
\end{tablenotes}
\end{threeparttable}
\end{table}

\end{document}
